%% kapa: exam=1 round=19 year=2008 period=1 form=A subject=accounting subject_ko=회계학 seq=22 number=102
%% source: Q-Net 감정평가사 기출문제 (cst00309) artlSeq=1001962 file=extracted/1008/1001962/2008년 제19회 감정평가사 1차시험 기출문제/1교시.hwp
%% conversion: kordoc 4.24.51 -> tools/gen1.py (auto); review=proofread-sonnet; flags=none
\begin{question}{102}
(주)도라지는 제품의 원가관리를 위하여 표준원가시스템을 적용하고 있다. 기준조업도는 월 1,000단위이고, 이를 기초로 계산된 단위당 제조간접원가 표준배부율은 ￦3,000이다. 1월 중 제품 생산량은 1,100단위이고, 이를 근거로 계산된 1월의 제조간접원가차이는 다음과 같다.

{\small\centering
\begin{tabular}{|ll|}
\hline
2분법 : & 제조간접원가 예산차이 : ￦150,000(유리) \\
 & 조업도차이 : ￦200,000(유리) \\
\hline
\end{tabular}\par}

2월 중 제품 생산량은 900단위이고, 실제 발생한 제조간접원가는 1월보다 ￦180,000이 적었다. 2월 중 제조간접원가 예산차이(유리)는 얼마인가?
\begin{choices}
  \item ￦70,000
  \item ￦90,000
  \item ￦110,000
  \item ￦130,000
  \item ￦170,000
\end{choices}
\end{question}
\answer{4}
